\documentclass[11pt]{article}
\usepackage[T1]{fontenc}
\usepackage{lmodern,amsmath,amssymb,amsthm,microtype}
\usepackage[margin=1in]{geometry}
\usepackage[colorlinks=true,linkcolor=blue,citecolor=blue,urlcolor=blue]{hyperref}
\newtheorem{theorem}{Theorem}
\newtheorem{lemma}[theorem]{Lemma}
\newtheorem{proposition}[theorem]{Proposition}
\newcommand{\R}{\mathbb R}
\newcommand{\ecc}{\operatorname{ecc}}
\title{Multiplicatively Weighted Centers of Median Graphs\\
in Quasilinear Time}
\author{Alec Kriebel\\\small Independent researcher\\
\small\href{https://orcid.org/0009-0001-9320-500X}{ORCID: 0009-0001-9320-500X}}
\date{October 14, 2026}
\begin{document}
\maketitle
\begin{abstract}
We record a short reduction from multiplicatively weighted center thresholds
to additively weighted eccentricities on a graph with bounded integral
distances. Combining it with the algorithm of Berg\'e, Ducoffe and Habib
and a known implicit optimization method gives a deterministic
$O(n\log^5(2n))$ algorithm for a center, the optimum radius and all centers
of every finite unit-edge median graph. This answers the finite algorithmic
form of Open Problem~1 in Ducoffe's Oberwolfach report, without a bound on
cube dimension. The contribution is the threshold observation and its
consequence; neither the additive algorithm nor the optimization mechanism
is new. We specify the exact comparison model and give explicit corrections
to two estimates in the public preprint proof of the imported algorithm.
\end{abstract}

\section{Problem and prior algorithms}
Let $G=(V,E)$ be a finite, nonempty, connected, simple graph with unit edge
lengths, $n=|V|$, and shortest-path distance $d$. For a finite nonnegative
real profile $w:V\to\R_{\geq0}$ define
\[
 r_w(x)=\max_{u\in V}w(u)d(x,u),\qquad
 \rho=\min_{x\in V}r_w(x),\qquad C_w=\{x:r_w(x)=\rho\}.
\]
A median graph is one in which every triple has a unique vertex belonging
to the three pairwise shortest-path intervals. Ducoffe~\cite[pp.~504--505]{OWR}
asks for an almost linear-time weighted-center algorithm on median graphs
of bounded cube dimension. The weights in this question multiply
distances; additive arrival weights define a different objective.

\begin{theorem}\label{thm:main}
Given a graph promised to be median, supplied by adjacency lists, and a
finite nonnegative real profile, one can compute $\rho$, a vertex in $C_w$
and the entire set $C_w$ deterministically in $O(n\log^5(2n))$ exact
operations. No cube-dimension bound or supplied embedding is required.
Operations on profile values only compare bounded integer multiples
$k w(u)$, with $0\leq k\leq n-1$.
\end{theorem}

The principal imported theorem is due to Berg\'e, Ducoffe and Habib
(BDH)~\cite{BDH}: for a nonnegative integer arrival profile $a$, all values
\begin{equation}\label{eq:additive}
 E_a(x)=\max_{u\in V}\bigl(d(x,u)+a(u)\bigr)
\end{equation}
on a median graph can be computed in $O(n\log^4(2n))$ time, including
preprocessing. The precise interface is on pp.~2--3 of the public
preprint~\cite{BDHarxiv}; its Algorithm~1 is on p.~22 and the overall
complexity proof is on pp.~23--24. Our calls have $0\leq a(u)\leq n-1$.
Appendix~\ref{app:repair} supplies consistent parameter and depth estimates
for two literal proof errors in that preprint version, preserving the
advertised theorem. Classical median structure and its preprocessing
algorithms remain credited dependencies~\cite{BDHarxiv,Median}.

Ducoffe~\cite[Lemma~2, p.~133:6]{D26} already proves that a $T(n,m)$-time
weighted-center decision algorithm yields optimization in
$O((T(n,m)+n)\log n)$ time. Its weighted-median mechanism is used below
in a slightly slower, self-contained comparison-only form. The full
preprint statement is Lemma~2.2 on pp.~5--6~\cite{D26arxiv}.
For cube-free median graphs the stronger bound $O(n\log^2 n)$ was already
proved by Chalopin et al.~\cite[Theorem~7.8, p.~28]{CCDDV}; our bound does
not improve that special case. Tree-product algorithms assume a stronger
representation promise~\cite{Trees,OWR}. The present note records the
short connection between these lines of work. A bounded primary-source
audit found no earlier explicit full-scope resolution, but does not
certify historical firstness or continued open status in October 2026.

\section{A threshold identity}
Set $D=n-1$. Every distance in a connected unit-edge graph is an integer
in $[0,D]$. For $R\geq0$ and $u\in V$ put
\begin{equation}\label{eq:budget}
 b_R(u)=\max\{k\in\{0,\ldots,D\}:k w(u)\leq R\},
 \qquad a_R(u)=D-b_R(u).
\end{equation}
Zero is an admissible index. If $w(u)=0$ then $b_R(u)=D$.
For positive $w(u)$ this is the capped floor of $R/w(u)$, but binary
search on the $D+1$ integer indices computes it using $O(\log(2n))$
comparisons, without a real division or unbounded-floor primitive.

\begin{lemma}\label{lem:threshold}
For every vertex $x$ and $R\geq0$,
\[
 r_w(x)\leq R\quad\Longleftrightarrow\quad E_{a_R}(x)\leq D.
\]
\end{lemma}
\begin{proof}
For every $u$, integrality and the distance bound give
\[
 w(u)d(x,u)\leq R
 \ \Longleftrightarrow\ d(x,u)\leq b_R(u)
 \ \Longleftrightarrow\ d(x,u)+a_R(u)\leq D.
\]
For a zero weight both conditions are vacuous. Take the conjunction
over $u$ and use~\eqref{eq:additive}.
\end{proof}

Thus one BDH call, followed by a scan, returns the \emph{entire} feasible
set $F_R=\{x:r_w(x)\leq R\}$ in $O(n\log^4(2n))$ operations.
The graph and its zero-weight vertices are retained. No subdivision,
Helly property, ball convexity or local-search theorem is required.
For $R<0$ the feasible set is empty. For $n=1$, or an all-zero profile,
the optimum is zero and every vertex is a center.

\section{Implicit optimization and complete output}
For any vertex $x$, its radius is an attained maximum and hence has the
form $k w(u)$ with $0\leq k\leq D$. Therefore $\rho$ belongs to the
multiset of $n$ sorted rows
\[
 S_u=(0,w(u),2w(u),\ldots,Dw(u)),\qquad u\in V.
\]
These $n^2$ entries are conceptual and are never materialized. Here is
the promised variant of the prior optimization mechanism~\cite{D26}.

Maintain an inclusive active index interval $[l_u,h_u]$ in each row,
initially $[0,D]$, and a feasible incumbent $B=D\max_u w(u)$, witnessed
by an arbitrary vertex. Handle $n=1$ and $B=0$ directly. Write
$L_u=h_u-l_u+1$ for a nonempty interval and let $M$ be their total length.
While $M>0$, take the lower middle value
$q_u=w(u)\lfloor(l_u+h_u)/2\rfloor$ of every nonempty row, with
selection weight $L_u$. Sort these values and choose a weighted median
$q$: the total selection weight on either weak side is at least $M/2$.
Query $F_q$ using Lemma~\ref{lem:threshold}.

If $F_q$ is nonempty, replace $B$ by $\min(B,q)$, saving a corresponding
feasible witness when it decreases, and delete all active entries
$\geq q$. Otherwise delete all active entries $\leq q$. Binary searches
within the row intervals implement either deletion exactly, including
ties and zero rows. Repeat until all intervals are empty.

\begin{lemma}\label{lem:search}
This process ends after $O(\log(2n))$ queries and returns $B=\rho$ and a
center. A final query at $B$ returns exactly $C_w$.
\end{lemma}
\begin{proof}
Maintain that $B$ is feasible and either $B=\rho$ or an active entry
equals $\rho$. Initially this holds. A feasible pivot has $\rho\leq q$:
if equality holds it is saved; if $\rho<q$, every optimum copy survives.
An infeasible pivot has $\rho>q$, so every optimum copy again survives.
This proves the invariant, and an empty active set forces $B=\rho$.
Its feasible witness is a center, and $F_B=C_w$.

For progress, a middle entry of a row of length $L_u$ has at least
$L_u/2$ entries on each weak side. If $q$ is infeasible, rows with
$q_u\leq q$ have total length at least $M/2$, and at least half of
those entries are deleted. If $q$ is feasible, use rows with $q_u\geq q$.
In either case at least $M/4$ entries are deleted, counting multiplicity.
This includes singleton rows and equal values. Since $M$ is an integer
and starts at $n^2$, only $O(\log(2n))$ iterations occur.
\end{proof}

Sorting and all bounded searches require $O(n\log(2n))$ work per
iteration and $O(n)$ row storage beyond the oracle. Thus the total
comparison overhead is $O(n\log^2(2n))$. The
$O(\log(2n))$ BDH calls and final complete-output call prove
Theorem~\ref{thm:main}. Reading the adjacency lists is absorbed: median
graphs have at most $n\log_2 n$ edges~\cite{BDHarxiv}.

\section{Model, reproducibility and status}
The claim concerns finite vertex centers in a graph promised to be
median. It includes zero weights, singleton support, exact equality
thresholds and multiple centers. It does not solve recognition, the
separate arbitrary-edge-length question, a continuous facility problem,
or an unspecified encoding of infinite median graphs. Nor does the
global threshold algorithm settle the separate unimodality question in
the workshop report~\cite{OWR}.

The optimum can be stored as one pair $(k,u)$ representing $k w(u)$.
All real comparisons are between such pairs. For rational input
$w(u)=p_u/q_u$ with positive denominators and components of at most $B$
bits, a comparison uses integer products $k p_u q_v$ and $j p_v q_u$
of $O(B+\log n)$ bits. There are $O(n\log^2(2n))$ original-weight
comparisons; their bit cost and input-reading cost must be paid.
The theorem is not an encoding-independent bit bound. The small
integer labels in the imported algorithm use $O(\log n)$ bits, including
its internal updates, as checked in Appendix~\ref{app:repair}.

The accompanying authored-only verification package includes the frozen
original candidate and historical audit, new independent proofs and
source comparisons, exact rational controls, the source-proof repair,
manuscript source, an offline replay driver and nested hash manifests.
Reference all-pairs and brute-structure routines verify finite cases;
they are expressly not a fast implementation or an experimental proof
of the asymptotic theorem. The universal argument uses the credited
classical structure and algorithms. Hashes establish file consistency,
not authenticity, priority or formal proof certification.

\paragraph{AI and review disclosure.}
AI tools were used extensively in solving, developing, writing,
reproducing and adversarially verifying this work. This is an
unrefereed preprint. Internal AI adversarial reviews are not human peer
review; the manuscript has not received human peer review or
proof-assistant certification. The priority assessment is bounded and
the limitations stated above are part of the result's presentation.

\appendix
\section{Explicit estimates for the imported recursion}\label{app:repair}
This appendix concerns the proof of Theorem~1 of
\emph{arXiv:2410.10235v1}~\cite{BDHarxiv}. It does not assert that the
final SODA typesetting contains the same errors; that full chapter was
not inspected. BDH's algorithm and structural discoveries remain theirs.
The complete step-by-step audit and exact controls are in the supplement.

That preprint declares $\log=\ln$ on p.~4, proves cube dimension
$d\leq\log_2 s$, then uses $\ell\leq\ln s$ for its wide-ladder size in
the duplicated-slice volume estimate on p.~24. Its separate displayed
product comparison in Lemma~16 on p.~13 has the wrong direction.
Both estimates admit the following consistent replacement.

For every recursive graph of order $s\geq3$ use
\[
 L(s)=\lceil\log_2 s\rceil,\qquad f(s)=2L(s).
\]
Compute $L(s)$ as the bit length of $s-1$. A cut is balanced exactly
when $f(s)$ times its smaller halfspace size is at least $s$.
Keep the original direct base cases $s\leq2$. The structural arguments
of BDH's Sections~3--4 apply with this choice: halfspaces and slices are
nonempty gated median graphs; the unique-majority median in an
unbalanced instance needs only $f(s)\geq2$. The wide ladder consists of
pairwise orthogonal incident classes, so $\ell\leq\log_2 s\leq L(s)$.
Each slice $S_i$ lies in a minority halfspace of size less than $s/f(s)$.
The two recursive calls on each slice consequently have total volume
\[
 2\sum_{i=0}^{\ell-1}|S_i|\leq \frac{2\ell s}{f(s)}\leq s.
\]
Balanced children partition $s$. Each child in either case has order at
most $s(1-1/f(s))$; in the unbalanced case its order is less than
$s/f(s)$. The gate, fiber-reweighting and merging identities depend on
gatedness, not this balance constant. Both original-weight and
fiber-weight slice calls are retained: a fiber weight need not dominate
the original value. There is no change of objective or input graph.

For completeness, put $\lambda=4/\ln2$. Since
$f(s)\leq4\log_2 s=\lambda\ln s$ for $s\geq3$, every nonterminal
path from a root of order $N$ satisfies
\[
 s_{i+1}\leq s_i\left(1-\frac1{\lambda\ln s_i}\right)
 \leq s_i\left(1-\frac1{\lambda\ln N}\right).
\]
Thus $s_h\leq N\exp(-h/(\lambda\ln N))$. At
$h=\lceil\lambda(\ln N)^2\rceil$ this is at most one, so the path has
already reached the $s\leq2$ base case. This proves depth $O(\log^2 N)$
directly, without the erroneous product comparison. The source's
nonrecursive work is $O(s\log^2(2s))$; per-level volume is at most $N$.
The total remains $O(N\log^4(2N))$.

For our initial offsets at most $N-1$, each fiber update adds at most
$N-1$ to a previous maximum. Along $O(\log^2 N)$ levels all internal
arrival labels are $O(N\log^2(2N))$, hence have $O(\log(2N))$ bits.
The cited linear-time median/Theta/halfspace preprocessing~\cite{Median}
is unchanged. It is not reconstructed as a fast implementation here.

\paragraph{Why the literal volume estimate needs correction.}
Take $Q_{30}$, attach $17{,}000{,}000{,}000$ leaves at each of its
30 unit-vector vertices and $488{,}926{,}258{,}176$ leaves at zero.
This is a median graph of order $N=10^{12}$: attaching pendant trees
at a vertex preserves medianity. Each coordinate minority has
$17{,}536{,}870{,}912$ vertices. A finite positive Taylor sum gives
$e^7>1000$, hence $\ln N<28$; multiplying each such minority by
56 gives $982{,}064{,}771{,}072<N$. All cuts are unbalanced for the
written natural-log threshold $2\ln N$. Give the cube antipode arrival
weight $N-1$ and every other vertex zero. It is uniquely farthest from
the majority median zero, with all 30 classes in its wide ladder.
Their minority union has $511{,}073{,}741{,}823$ vertices; its two
slice copies therefore total $1{,}022{,}147{,}483{,}646>N$.
This succinct exact example refutes the intermediate volume estimate,
not correctness or the final time bound after the stated repair.

\begin{thebibliography}{9}
\raggedright
\bibitem{OWR} G. Ducoffe, Radius functions in median graphs, in
\emph{Median Geometry and Applications}, Oberwolfach Reports 8/2026,
504--507. \href{https://doi.org/10.4171/OWR/2026/8}{doi:10.4171/OWR/2026/8}.
\bibitem{BDH} P. Berg\'e, G. Ducoffe and M. Habib,
Quasilinear-time eccentricities computation, and more, on median graphs,
\emph{SODA 2025}, 1679--1704.
\href{https://doi.org/10.1137/1.9781611978322.52}{doi:10.1137/1.9781611978322.52}.
\bibitem{BDHarxiv} P. Berg\'e, G. Ducoffe and M. Habib, same title,
public preprint, October 14, 2024.
\href{https://arxiv.org/abs/2410.10235v1}{arXiv:2410.10235v1}.
\bibitem{D26} G. Ducoffe, Beyond Trees: The Weighted Center Problem on
Gromov Hyperbolic Graphs, \emph{ESA 2026}, LIPIcs 388, 133:1--133:19.
\href{https://doi.org/10.4230/LIPIcs.ESA.2026.133}{doi:10.4230/LIPIcs.ESA.2026.133}.
\bibitem{D26arxiv} G. Ducoffe, same title, full preprint, July 5, 2026.
\href{https://arxiv.org/abs/2607.04287v1}{arXiv:2607.04287v1}.
\bibitem{CCDDV} J. Chalopin, V. Chepoi, F. F. Dragan, G. Ducoffe and
Y. Vax\`es, On $G^p$-unimodality of radius functions in graphs:
structure and algorithms, March 19, 2025.
\href{https://arxiv.org/abs/2503.15011v1}{arXiv:2503.15011v1}.
\bibitem{Trees} G. Ducoffe, Isometric Embeddings in Trees and Their Use
in Distance Problems, \emph{MFCS 2021}, LIPIcs 202, 43:1--43:16.
\href{https://doi.org/10.4230/LIPIcs.MFCS.2021.43}{doi:10.4230/LIPIcs.MFCS.2021.43}.
\bibitem{Median} L. B\'en\'eteau, J. Chalopin, V. Chepoi and Y. Vax\`es,
Medians in median graphs and their cube complexes in linear time,
\emph{J. Comput. Syst. Sci.} 126 (2022), 80--105.
\href{https://doi.org/10.1016/j.jcss.2022.01.001}{doi:10.1016/j.jcss.2022.01.001}.
\end{thebibliography}
\end{document}
